\chapter{Diseño factorial (si aplica)}
\label{anexo:e}

El diseño factorial completo \textbf{no aplica} al presente plan de tesis por las siguientes
razones:

\begin{itemize}
  \item El estudio es de tipo \textbf{investigación tecnológica / design science}: el producto del
        plan es el diseño del procedimiento y del modelo S/E/C, no un experimento factorial con
        población y muestra al estilo cuasi-experimental.
  \item Los \textbf{factores} (esquema de representación, técnica de matching, pesos
        $\alpha/\beta/\gamma$) se documentan como \textbf{decisiones de diseño} con estados
        alternativos posibles, pero su evaluación comparativa sistemática corresponde a una etapa
        de validación empírica explícitamente fuera de alcance (§1.6).
  \item La matriz de operación de variables (Anexo~D) y la matriz de consistencia (Anexo~C)
        sustituyen al diseño factorial para documentar el diseño del procedimiento.
\end{itemize}

Si en una fase posterior se requiriese un diseño experimental, los factores candidatos serían:
(1)~representación intermedia, (2)~técnica semántica, (3)~estrategia estructural y
(4)~ponderación $\alpha,\beta,\gamma$; ello excede el alcance del presente documento.
